% =====================================================================
%  1bat-ud1-3.tex  —  HORA 3: El bo social elèctric (descoberta dels intervals)
% =====================================================================
\horatitol{Sessió 3 — El bo social elèctric: descoberta dels intervals}%
{Un cas real: decidir qui té dret a un ajut a partir d'un llindar de renda, abans de cap fórmula.}

\subsection{Idea clau}

Avui treballem com a investigadors, en grup i a les pissarres. Tenim un cas
\textbf{real}: el \textbf{bo social elèctric}, un descompte a la factura de la llum
per a famílies amb pocs ingressos. El dret a l'ajut depèn d'un \textbf{llindar de
renda} fixat a partir de l'\textbf{IPREM} (un indicador oficial que es fa servir per
fixar molts ajuts). El repte: decidir qui hi té dret \emph{sense} que us doni cap
fórmula. L'heu de trobar vosaltres.

\begin{idea}
\textbf{Context} (xifres simplificades per treballar; cal actualitzar-les al curs vigent):
\begin{itemize}[leftmargin=1.4em, itemsep=2pt]
  \item L'IPREM anual de referència és, posem, $\num{8400}\eur$.
  \item Una família \textbf{sense} circumstàncies especials té dret al bo si la seva
        renda anual \textbf{no supera} $1,5$ vegades l'IPREM.
  \item Una família \textbf{nombrosa} té un llindar més alt: fins a $2$ vegades l'IPREM.
\end{itemize}
\textbf{La pregunta del dia:} com decidiu, per a cada família, si hi té dret? I com
ho escriuríeu \emph{amb símbols}, sense paraules com «fins a» o «com a màxim»?
\end{idea}

\subsection{Exemples}

\begin{exemple}[calcular el llindar]
Primer cal saber \emph{quin} és el llindar de cada cas. Amb un IPREM de $\num{8400}\eur$:
\begin{itemize}[leftmargin=1.4em, itemsep=2pt]
  \item Família ordinària: $1,5\times\num{8400}=\num{12600}\eur$. Té dret si la renda
        \textbf{no supera} $\num{12600}\eur$.
  \item Família nombrosa: $2\times\num{8400}=\num{16800}\eur$. Té dret si la renda no
        supera $\num{16800}\eur$.
\end{itemize}
«No superar» vol dir que el límit \emph{sí} que hi entra: una renda \emph{just} igual
al llindar encara dóna dret.
\end{exemple}

\begin{exemple}[decidir un perfil]
La família Garcia (ordinària) té una renda anual de $\num{11200}\eur$. Com que
$\num{11200}\le\num{12600}$, \textbf{sí} que té dret al bo social. En canvi, la
família Soler (ordinària) amb $\num{13000}\eur$ \textbf{no} hi té dret, perquè
$\num{13000}>\num{12600}$ (supera el llindar).
\end{exemple}

\subsection{Exercicis resolts}

\begin{exresolt}[decidir amb el llindar ordinari]
Amb el llindar ordinari de $\num{12600}\eur$, digues si tenen dret al bo aquestes
famílies i escriu la condició amb símbols: (a) renda $\num{9800}\eur$; (b) renda
$\num{12600}\eur$; (c) renda $\num{14500}\eur$.
\solucio
La condició per tenir-hi dret és $\text{renda}\le\num{12600}$.
\begin{enumerate}[label=\alph*), leftmargin=1.6em]
  \item $\num{9800}\le\num{12600}$: \textbf{sí}.
  \item $\num{12600}\le\num{12600}$: \textbf{sí} (just al límit, que hi entra).
  \item $\num{14500}\le\num{12600}$ és fals: \textbf{no}.
\end{enumerate}
\resposta{(a) sí; (b) sí (el límit hi entra); (c) no.}
\end{exresolt}

\begin{exresolt}[el llindar canvia amb les circumstàncies]
La família Mbeki és \textbf{nombrosa} i té una renda de $\num{15500}\eur$. Tindria
dret al bo? I si \emph{no} fos nombrosa?
\solucio
Com a família nombrosa, el seu llindar és $2\times\num{8400}=\num{16800}\eur$. Com
que $\num{15500}\le\num{16800}$, \textbf{sí} que hi té dret.

Si \emph{no} fos nombrosa, el llindar seria $\num{12600}\eur$, i com que
$\num{15500}>\num{12600}$, \textbf{no} hi tindria dret. La mateixa renda dóna un
resultat diferent segons les circumstàncies: el llindar és el que canvia.
\resposta{Com a nombrosa, sí (llindar $\num{16800}$); si no ho fos, no (llindar $\num{12600}$).}
\end{exresolt}

\subsection{Exercicis per fer}

\noindent\textit{Treballeu en grup. Useu un IPREM de $\num{8400}\eur$; llindar
ordinari $1,5\times$IPREM, llindar de família nombrosa $2\times$IPREM.}

\begin{exercicis}
\begin{exlist}
  \item Calcula els dos llindars en euros: el de família ordinària ($1,5\times$IPREM)
        i el de família nombrosa ($2\times$IPREM).
  \item Decideix si tenen dret al bo (família ordinària, llindar $\num{12600}\eur$):
        \begin{enumerate}[label=\alph*), leftmargin=1.6em, topsep=2pt]
          \item renda $\num{8000}\eur$ \hfill \item renda $\num{12600}\eur$ \hfill
          \item renda $\num{12601}\eur$
        \end{enumerate}
  \item Escriu \textbf{amb símbols} (sense paraules) la condició per tenir dret al bo
        en cada cas:
        \begin{enumerate}[label=\alph*), leftmargin=1.6em, topsep=2pt]
          \item família ordinària
          \item família nombrosa
        \end{enumerate}
  \item La família Costa té $\num{16800}\eur$ de renda. Tindria dret si fos nombrosa?
        I si no ho fos? Justifica-ho amb una desigualtat en cada cas.
  \item \textbf{(Frontera exacta)} Dues famílies ordinàries tenen rendes de
        $\num{12600}\eur$ i $\num{12600,50}\eur$. Quina té dret al bo i quina no? Què
        ens diu això sobre la importància d'on posem exactament el límit?
  \item \textbf{(Volta de pissarres)} El teu grup ha escrit la condició d'una manera;
        un altre grup, potser, d'una altra. Anota \textbf{dues maneres diferents}
        d'expressar «la renda no supera $\num{12600}\eur$» que hagin sortit a classe
        (amb paraules o amb símbols).
\end{exlist}
\end{exercicis}

% ---------------------------------------------------------------------
\fullsolucions{Sessió 3}
\begin{solucions}
\begin{sollist}
  \item Ordinària: $1,5\times\num{8400}=\num{12600}\eur$; nombrosa:
        $2\times\num{8400}=\num{16800}\eur$.
  \item (a) sí ($\num{8000}\le\num{12600}$); \quad (b) sí (just al límit); \quad
        (c) no ($\num{12601}>\num{12600}$).
  \item (a) $\text{renda}\le\num{12600}$; \quad (b) $\text{renda}\le\num{16800}$.
  \item Com a nombrosa: $\num{16800}\le\num{16800}$, sí (just al límit). Si no ho fos:
        $\num{16800}>\num{12600}$, no.
  \item La de $\num{12600}\eur$ té dret ($\le$ el límit); la de $\num{12600,50}\eur$ no
        ($>$ el límit). Per pocs cèntims es perd el dret: per això cal definir
        \emph{exactament} el llindar i si el límit hi entra o no.
  \item Resposta oberta. Exemples: «fins a $\num{12600}\eur$», «no més de
        $\num{12600}\eur$», «renda $\le\num{12600}$», «$\le 1,5\times$IPREM».
\end{sollist}
\end{solucions}
